\clearpage
\section{연습문제}
\label{sec:radicals-solvable-exercises}

문제는 A, B, C의 세 단계로 나뉜다. 별표 문제는 다음 두 절에 상세 풀이가 있다.

\subsection*{A. 계산과 구체적 판정}

\begin{exercise}[*]
$S_3$과 $S_4$의 도함열을 계산하고 각각의 도함길이를 구하라. $S_4$의 경우
\[
  S_4\trianglerighteq A_4\trianglerighteq V_4
  \trianglerighteq1
\]
과 도함열의 관계를 설명하라.
\end{exercise}

\begin{exercise}
다음 군이 가해인지 판정하고, 가해인 경우 소수차 순환 인자군을 갖는 정규열을 하나 제시하라.
\[
  C_{12},\qquad D_5,\qquad A_4,
  \qquad F_{20}=C_5\rtimes C_4,
  \qquad A_5.
\]
\end{exercise}

\begin{exercise}[*]
다항식 $X^3-2$의 분해체
\[
  L=\Q(\zeta_3,\sqrt[3]{2})
\]
에 대해 다음을 수행하라.
\begin{enumerate}[label=(\alph*)]
  \item $\Gal(L/\Q)\cong S_3$임을 보여라.
  \item 정규열 $S_3\trianglerighteq A_3\trianglerighteq1$에 대응하는 중간체를 구하라.
  \item $L/\Q$를 포함하는 근호탑을 적고 각 단계의 종류를 설명하라.
\end{enumerate}
\end{exercise}

\begin{exercise}
분리다항식의 갈루아군이 다음과 같다고 알려졌을 때 방정식이 근호로 풀리는지 판정하라.
\[
  C_7,\quad V_4,\quad D_4,\quad A_4,\quad S_4,
  \quad D_5,\quad F_{20},\quad A_5,\quad S_5.
\]
\end{exercise}

\begin{exercise}[*]
$f=X^5-4X+2$에 대해 다음을 증명하라.
\begin{enumerate}[label=(\alph*)]
  \item $f$는 $\Q$ 위에서 기약이다.
  \item $f$는 정확히 세 개의 실근과 두 개의 비실근을 갖는다.
  \item $f(x)=0$은 근호로 풀리지 않는다.
  \item 갈루아군이 $S_5$임을 보여라.
\end{enumerate}
\end{exercise}

\begin{exercise}
$f=X^5-2$의 분해체를 구하고 갈루아군이
\[
  C_5\rtimes C_4
\]
의 부분군임을 보여라. 이 다항식이 오차식임에도 근호로 풀리는 이유를 설명하라.
\end{exercise}

\begin{exercise}[*]
삼차식
\[
  X^3-3X+1=0
\]
에 대해 판별식과 갈루아군을 구하라. Cardano 공식에서 $u^3,v^3$를 계산하고, $uv=1$이 되도록 세제곱근을 택해 세 근을 적어라.
\end{exercise}

\begin{exercise}
삼차식
\[
  X^3+3X-4=0
\]
을 유리근을 이용해 풀고 Cardano 공식으로 다시 풀어라. Cardano 식에서 나타나는 근호들이 최종적으로 어떻게 단순화되는지 설명하라.
\end{exercise}

\begin{exercise}[*]
사차식
\[
  X^4-5X^2+6=0
\]
의 삼차분해식을 계산하고 완전히 인수분해하라. 그 세 근 $z_1,z_2,z_3$에서 $\sqrt{-z_i}$를 택하여 원래 네 근을 복원하라.
\end{exercise}

\begin{exercise}
사차식 $X^4-2$의 삼차분해식을 구하고, 분해체
\[
  \Q(\sqrt[4]{2},i)
\]
의 근호탑과 갈루아군 $D_4$의 정규열을 대응시켜라.
\end{exercise}

\subsection*{B. 핵심 정리의 증명}

\begin{exercise}[*]
교환자부분군 $[G,G]$가 정규부분군이고 $G/[G,G]$가 아벨군임을 증명하라. 또한 $G/N$이 아벨인 모든 정규부분군 $N$에 대해
\[
  [G,G]\subseteq N
\]
임을 증명하라.
\end{exercise}

\begin{exercise}
가해군의 부분군과 몫군이 가해군임을 도함열로 증명하라. $N\triangleleft G$이고 $N,G/N$이 가해이면 $G$도 가해임을 보여라.
\end{exercise}

\begin{exercise}[*]
군 $G$가 가해일 필요충분조건이 어떤 $r$에 대해
\[
  G^{(r)}=1
\]
인 것임을 증명하라. 정규열과 도함열 사이의 포함관계를 명시하라.
\end{exercise}

\begin{exercise}
유한 $p$-군의 중심이 비자명하다는 사실을 사용하여 모든 유한 $p$-군이 가해군임을 군의 위수에 대한 귀납법으로 증명하라.
\end{exercise}

\begin{exercise}[*]
유한 가해군이 인자군이 모두 소수차 순환군인 정규열을 가짐을 증명하라. 유한성 가정이 세분과정의 종료에 어떻게 사용되는지 설명하라.
\end{exercise}

\begin{exercise}
다음을 증명하라.
\[
  [S_n,S_n]=A_n\qquad(n\ge2),
\]
\[
  [A_n,A_n]=A_n\qquad(n\ge5).
\]
이를 이용해 $S_n$이 $n\ge5$에서 가해군이 아님을 보여라.
\end{exercise}

\begin{exercise}[*]
유한확장 $L/K$가 가해 또는 근호로 풀릴 때, 임의의 기준체 확대 $F/K$에 대해 합성체 $FL/F$도 같은 성질을 가짐을 증명하라. 이어 유한확장탑
\[
  K\subseteq L\subseteq M
\]
에서 가해성과 근호해법이 위아래로 보존됨을 설명하라.
\end{exercise}

\begin{exercise}
근호 단계의 세 유형이 모두 가해확장임을 증명하라. 특히 $\alpha^n=a\in K$인 단계에서는 원시 $n$차 단위근을 첨가한 정상폐포의 갈루아군이 순환군에 의한 아벨군의 확대임을 보여라.
\end{exercise}

\begin{exercise}[*]
유한확장 $L/K$가 가해이면 근호로 풀림을 다음 순서로 증명하라.
\begin{enumerate}[label=(\alph*)]
  \item $L$을 가해 갈루아확장 $E/K$ 안에 넣어라.
  \item $[E:K]$의 표수와 다른 소인수들에 대한 단위근을 첨가하라.
  \item 갈루아군의 정규열을 소수차 순환 인자군들로 세분하라.
  \item 갈루아 대응과 Kummer·아르틴--슈라이어 정리를 적용하라.
\end{enumerate}
역방향도 근호 단계의 정상폐포를 이용해 증명하라.
\end{exercise}

\begin{exercise}
차수 $4$ 이하인 모든 유한 분리확장이 근호로 풀림을 원시원 정리와 $S_4$의 가해성으로 증명하라.
\end{exercise}

\subsection*{C. 소수차, 분해식과 구조적 응용}

\begin{exercise}[*]
$p$가 소수이고 $G\le S_p$가 가해이며 추이적이라고 하자. $G$가 유일한 정규 $p$-부분군을 가짐을 보이고, 점 집합을 $\F_p$와 동일시하여
\[
  G\le\operatorname{AGL}(1,p)
\]
임을 증명하라.
\end{exercise}

\begin{exercise}
가해 추이적 부분군 $G\le S_p$의 비항등원이 고정점을 많아야 하나만 가짐을 증명하라. 이를 이용해 가해인 기약 소수차다항식의 분해체가 임의의 서로 다른 두 근으로 생성됨을 보여라.
\end{exercise}

\begin{exercise}[*]
다항식
\[
  f=X^7-6X+3\in\Q[X]
\]
에 대해 다음을 수행하라.
\begin{enumerate}[label=(\alph*)]
  \item $f$가 기약임을 보여라.
  \item 도함수를 이용해 실근이 정확히 세 개임을 보여라.
  \item 방정식이 근호로 풀리지 않음을 결론 내리라.
\end{enumerate}
\end{exercise}

\begin{exercise}
$p$가 소수이고 $G\le S_p$가 가해이며 추이적이면
\[
  |G|\mid p(p-1)
\]
임을 증명하라. 가능한 $p$-Sylow 부분군과 점안정자의 구조를 설명하라.
\end{exercise}

\begin{exercise}[*]
기약 오차식의 갈루아군이 $D_5$라고 하자.
\begin{enumerate}[label=(\alph*)]
  \item $D_5\trianglerighteq C_5\trianglerighteq1$이 가해 정규열임을 보여라.
  \item 갈루아 대응으로 대응하는 체의 탑을 적어라.
  \item 필요한 단위근을 첨가한 뒤 각 단계를 어떤 근호로 표현할 수 있는지 설명하라.
\end{enumerate}
\end{exercise}

\begin{exercise}
$F_{20}=C_5\rtimes C_4$의 교환자부분군과 도함열을 구하라. $F_{20}$이 가해이지만 아벨이 아님을 확인하라.
\end{exercise}

\begin{exercise}[*]
$f\in K[X]$가 기약 오차다항식이고 갈루아군이 $A_5$라고 하자.
\begin{enumerate}[label=(\alph*)]
  \item 판별식이 $K$의 제곱임을 보여라.
  \item $A_5$가 완전군이라는 사실에서 $f$가 근호로 풀리지 않음을 증명하라.
  \item 판별식이 제곱이라는 조건만으로는 근호해법을 보장하지 않는 이유를 설명하라.
\end{enumerate}
\end{exercise}

\begin{exercise}
$X^5-2$와 일반 오차방정식을 비교하라. 두 갈루아군의 정규열을 비교하여 한쪽은 근호로 풀리고 다른 쪽은 풀리지 않는 구조적 이유를 설명하라.
\end{exercise}

\begin{exercise}
$X^3+pX+q$의 근 $x_1,x_2,x_3$과 원시 세제곱근 $\zeta$에 대해
\[
  R=x_1+\zeta x_2+\zeta^2x_3,
  \qquad
  S=x_1+\zeta^2x_2+\zeta x_3
\]
라 하자. 직접 전개하여
\[
  RS=-3p
\]
와 $R^3,S^3$의 Cardano 항등식을 증명하라.
\end{exercise}

\begin{exercise}[*]
우울형 사차식
\[
  X^4+pX^2+qX+r
\]
의 네 근에서 세 양 $z_1,z_2,z_3$을 정의하고 다음을 증명하라.
\begin{enumerate}[label=(\alph*)]
  \item $z_i$들은
  \[
    Z^3-2pZ^2+(p^2-4r)Z+q^2
  \]
  의 근이다.
  \item $a^2=-z_1$, $b^2=-z_2$, $c^2=-z_3$, $abc=-q$이면 네 근이
  \[
    \frac12(\pm a\pm b\pm c)
  \]
  의 적절한 네 부호조합으로 복원된다.
  \item 이 과정이 $S_4/V_4$와 $V_4$의 두 층을 어떻게 반영하는지 설명하라.
\end{enumerate}
\end{exercise}
